\documentclass[11pt]{article}

\usepackage[margin=1in]{geometry}
\usepackage[T1]{fontenc}
\usepackage[utf8]{inputenc}
\usepackage{amsmath,amssymb,amsthm,mathtools}
\usepackage{array}
\usepackage{booktabs}
\usepackage{enumitem}
\usepackage{graphicx}
\usepackage{longtable}
\usepackage[expansion=false]{microtype}
\usepackage{tikz}
\usetikzlibrary{arrows.meta,positioning,shapes.geometric}
\usepackage{hyperref}
\usepackage[numbers,sort&compress]{natbib}
\usepackage{cleveref}
\emergencystretch=2em

\newcolumntype{L}[1]{>{\raggedright\arraybackslash\sloppy}p{#1}}
\setlength{\tabcolsep}{5pt}
\renewcommand{\arraystretch}{1.12}

\hypersetup{
  colorlinks=true,
  linkcolor=blue,
  citecolor=blue,
  urlcolor=blue
}

\newcommand{\E}{\mathbb{E}}
\newcommand{\PPNR}{\mathrm{PPNR}}
\newcommand{\EL}{\mathrm{EL}}
\newcommand{\NPV}{\mathrm{NPV}}
\newcommand{\1}{\mathbf{1}}
\newcommand{\code}[1]{\nolinkurl{#1}}
\DeclareMathOperator*{\argmax}{arg\,max}
\DeclareMathOperator*{\argmin}{arg\,min}
\DeclareMathOperator{\softmax}{softmax}
\DeclareMathOperator{\diag}{diag}

\theoremstyle{definition}
\newtheorem{definition}{Definition}[section]
\newtheorem{assumption}[definition]{Assumption}
\newtheorem{proposition}[definition]{Proposition}
\newtheorem{remark}[definition]{Remark}

\title{A Full Bayesian State-Space and Belief-State Control Specification\\
for Credit-Card Customer NPV Decisioning}
\author{Chak Wong}
\date{July 2, 2026}

\begin{document}
\maketitle

\begin{abstract}
This note gives a full stochastic-process assembly of the credit-card customer
NPV proposal into a constrained Bayesian state-space model with belief-state
control. The goal is not to replace the proposal's valuation semantics. It is
to make the dynamic object explicit enough that forward simulation under a
fixed declared policy is mathematically well defined and later offline policy
evaluation or constrained challenger optimization can be discussed without
handwaving. The note reuses the current proposal's valuation object, wallet
accounting, module decomposition, path-engine semantics, and policy-value
conventions, then adds the missing layers required for a committee-level
stochastic specification: timing and filtration, explicit latent and observed
state blocks, module-level controlled transition equations, observation
equations, Bayesian filtering recursions, constrained Bellman equations, and an
identification / governance boundary. The resulting object is a lifecycle
belief-state control model, not a claim that the bank is already ready for
production reinforcement learning. The literature supports the sequential
control framing, partial observability, state-space filtering, and offline
policy-evaluation toolkit
\citep{puterman1994mdp,murphy2003dynamic,smallwood1973pomdp,kaelbling1998pomdp,
trench2003bankone,durbinkoopman2012ssm,westharrison1997bfdm,doucet2001smc,
jiangli2016dr,huwager2021pomdpope,laroche2019spibb,khraishiokhrati2022offlinecredit};
it does not identify this issuer's exact latent state, reward ledger,
constraint set, or safe rollout threshold.
\end{abstract}

\section{Introduction and contribution boundary}

The credit-card NPV proposal already contains the essential economic pieces of a
dynamic decision system: a versioned valuation object, a modular state
structure, wallet accounting, action-dependent terms, a path engine, a
cash-flow compiler, and policy-value notation. What it does not yet present is a
single stochastic process that assembles those pieces into one mathematically
explicit control object. Without that object the proposal can say that future
policy matters, but it cannot fully specify what is being filtered, what is
transitioning, what is observed with noise or lag, or exactly which state a
policy should condition on.

That gap matters because many committee-level objections are structural rather
than stylistic. If the proposal names outside wallet, movable wallet, repayment
propensity, rewards response, attrition, PD/LGD/EAD, and macro sensitivity as
important drivers, then a mathematically serious note must say where those
objects live in the state, how they evolve, how they are observed, and how they
enter the reward and feasible action set. A high-level POMDP discussion is not
enough.

This note therefore contributes a full stochastic-process specification. It
makes five bounded claims.
\begin{enumerate}[label=\arabic*.]
  \item Credit-card customer value can be represented as a sequential control
        problem, with discounted incremental NPV as the policy-value object
        \citep{trench2003bankone,puterman1994mdp,murphy2003dynamic}.
  \item A fully observed MDP is generally an approximation because important
        economic drivers are latent or only partially observed; a belief-state
        or POMDP formulation is therefore the correct umbrella language
        \citep{smallwood1973pomdp,kaelbling1998pomdp}.
  \item The current proposal's macro, spend, balance, payment, loss, reward,
        attrition, and relationship modules can be assembled into one controlled
        state-space system.
  \item Bayesian filtering supplies the correct updating object: the policy
        should depend on the posterior distribution of the latent state, not
        only on a point estimate
        \citep{durbinkoopman2012ssm,westharrison1997bfdm,doucet2001smc}.
  \item Offline policy evaluation and safe-improvement tools are relevant only
        after the dynamic system, logged behavior policy, and support conditions
        are explicit
        \citep{jiangli2016dr,huwager2021pomdpope,laroche2019spibb,
        khraishiokhrati2022offlinecredit}.
\end{enumerate}

The note does \emph{not} claim that the literature already identifies this
issuer's exact latent state vector, its competitor-wallet dynamics, its future
policy bundle, or its live deployment threshold. Those remain bank-specific
empirical and governance questions.

\section{Notation and timing contract}

This note reuses the main proposal's global notation. Unit $i$ is a customer,
account, prospect, application, or account-action unit. Time $t$ is monthly.
Action is $a$. The decision-specific baseline is $a_0(d)$. Decision context is
$d$. Scenario and valuation bundle is $s$. Downstream account-management policy
bundle is $\pi^{down}$. Governance and hard admissibility are represented by
$\pi^{gov}$. The decision-time information set is $\mathcal{I}_{id}$ or
$X_{it}^{d}$.

The current proposal's one-period incremental cash-flow primitive and NPV
functional are retained:
\begin{align}
  \Delta CF_{i,t+h}(a,\pi;s)
  &=
  \Delta \PPNR_{i,t+h}(a,\pi;s)
  - \Delta \EL_{i,t+h}(a,\pi;s)
  - \Delta Kchg_{i,t+h}(a,\pi;s)
  \nonumber\\
  &\quad
  - \Delta Tax_{i,t+h}(a,\pi;s)
  + \Delta RelValue_{i,t+h}(a,\pi;s),
  \label{eq:cf-primitive}\\
  \Delta \NPV_i(a;d,s,\pi)
  &=
  -C_i^{\mathrm{acq}}(a)
  +
  \E\!\left[
    \sum_{h=0}^{H}\delta_h\Delta CF_{i,t+h}(a,\pi;s)
    + \delta_H\Delta TV_{i,t+H}(a,\pi;s)
    \mid X_{it}^{d}
  \right].
  \label{eq:npv-functional}
\end{align}

This note distinguishes three policy objects. The symbol $\pi^{down}$ denotes
the downstream operating policy bundle assumed in valuation. The symbol $\mu$
denotes the historical logged behavior policy that generated observed actions.
The symbol $\varphi$ denotes a candidate or challenger rule to be simulated,
evaluated, or eventually optimized. When the older proposal notation uses
$\pi$ generically, the reader should check whether it refers to $\pi^{down}$,
$\mu$, $\varphi$, or the governance bundle $\pi^{gov}$.

\subsection{Intra-month timing}

A committee-level dynamic model needs an explicit within-month sequence. Let the
month-$t$ cycle be:
\begin{enumerate}[label=\textbf{T\arabic*.},leftmargin=*]
  \item At the opening of month $t$, the bank has observed history
        $H_{it}^{-}$ and a posterior belief $b_{it}^{-}$ over the latent state.
  \item Macro or scenario inputs available at decision time are incorporated;
        lagged or not-yet-released macro variables remain unavailable.
  \item The bank constructs the decision feature state $O_{it}$ and chooses an
        admissible action
        \[
          a_{it}\in\mathcal{A}_t(O_{it},b_{it}^{-};d,s,\pi^{gov}).
        \]
  \item Conditional on $a_{it}$ and the current latent/full state, within-month
        customer behavior realizes: activation or use, category spend,
        cash-advance or balance-transfer use, balances, payments, delinquency
        migration, rewards accrual, attrition or retention outcomes, and any
        collections response.
  \item Raw signals $Y_{i,t+1}$ are observed with the data lags, censoring, or
        partial observability implied by the source system.
  \item The posterior is updated to $b_{i,t+1}^{-}$ using Bayes' rule, and the
        month-$t$ reward is compiled from the realized or expected incremental
        cash-flow mapping.
\end{enumerate}

The relevant filtration should be written in two steps. The pre-action
information set is
\begin{equation}
  \mathcal{F}^{-}_{it}
  =
  \sigma(H_{it}^{-},O_{it},b_{it}^{-}),
  \label{eq:pre-action-filtration}
\end{equation}
and the action must be measurable with respect to it:
\begin{equation}
  a_{it}=\varphi_t(O_{it},b_{it}^{-},d,s,\pi^{down})
  \quad\text{or}\quad
  a_{it}\sim\varphi_t(\cdot\mid O_{it},b_{it}^{-},d,s,\pi^{down}).
  \label{eq:policy-measurability}
\end{equation}
After the action and raw signals are realized, the post-observation information
set is
\begin{equation}
  \mathcal{F}^{+}_{i,t+1}
  =
  \sigma(\mathcal{F}^{-}_{it},a_{it},Y_{i,t+1},O_{i,t+1}).
  \label{eq:post-action-filtration}
\end{equation}

\section{State decomposition}

\subsection{Observed state}

Let the observed decision state be
\begin{equation}
  O_{it}
  =
  (G_{it},B_{it},L_{it},P_{it},D_{it},Rwd_{it},Rel^{obs}_{it},M^{obs}_{it},C_{it}),
  \label{eq:observed-state}
\end{equation}
where:
\begin{itemize}[leftmargin=*]
  \item $G_{it}$ is lifecycle regime,
  \item $B_{it}$ is the balance-subledger vector,
  \item $L_{it}$ is line and term state,
  \item $P_{it}$ is observed payment behavior,
  \item $D_{it}$ is observed delinquency/default/cure status,
  \item $Rwd_{it}$ is observed reward and promotion state,
  \item $Rel^{obs}_{it}$ is observed relationship state,
  \item $M^{obs}_{it}$ is the observed macro signal vector available at decision
        time, and
  \item $C_{it}$ is decision-context metadata such as channel, offer family,
        feature-cut timestamp, and eligibility indicators.
\end{itemize}

The balance-subledger vector should be explicit:
\begin{equation}
  B_{it}
  =
  (B^{pur}_{it},B^{cash}_{it},B^{BT}_{it},B^{fee}_{it},B^{int}_{it},B^{chg}_{it}),
  \label{eq:balance-subledger}
\end{equation}
so that principal, non-principal, promo, and charge-related balances are not
collapsed into one scalar.

The lifecycle regime should also be explicit:
\begin{equation}
  \begin{split}
  G_{it}\in\{&
    \text{prospect},\text{targeted},\text{applicant},
    \text{approved},\text{booked},\text{activated},
    \text{active-transactor},\\
    &\text{active-revolver},\text{promo},\text{dormant},
    \text{retained},\text{converted},\text{delinquent},
    \text{collections},\\
    &\text{closed},\text{charged-off}\}.
  \end{split}
  \label{eq:lifecycle-regime}
\end{equation}

\subsection{Latent state}

Let the latent economic state be
\begin{equation}
  U_{it}
  =
  (M^{lat}_{it},W^{out}_{it},\Theta^{route}_{it},\Theta^{repay}_{it},
   \Theta^{attr}_{it},\Theta^{promo}_{it},Rel^{lat}_{it}),
  \label{eq:latent-state}
\end{equation}
where:
\begin{itemize}[leftmargin=*]
  \item $M^{lat}_{it}$ is a latent macro or local economic regime, if the bank
        chooses to smooth or nowcast macro conditions beyond raw releases,
  \item $W^{out}_{it}$ is outside wallet or outside balance opportunity,
  \item $\Theta^{route}_{it}$ captures latent routing / card-on-file / issuer-
        preference frictions,
  \item $\Theta^{repay}_{it}$ captures latent repayment, liquidity, and
        distress state not fully summarized by observed balances and payments,
  \item $\Theta^{attr}_{it}$ captures latent attrition / dormancy / retention
        propensity,
  \item $\Theta^{promo}_{it}$ captures latent promotion-response and decay
        state, and
  \item $Rel^{lat}_{it}$ captures latent relationship spillover state.
\end{itemize}

The full economic state is
\begin{equation}
  X_{it} = (O_{it},U_{it}).
  \label{eq:full-state}
\end{equation}

\subsection{Belief state}

Because $U_{it}$ is not observed, the control object is the posterior
belief state
\begin{equation}
  b^{-}_{it}(u)
  =
  p(U_{it}=u\mid H_{it}^{-},s,\pi^{gov}).
  \label{eq:belief-state}
\end{equation}
The policy may use a finite-dimensional approximation to $b_{it}$ in
implementation, but the conceptual sufficient statistic is the posterior
distribution itself.
For readability, later equations usually write $b_{it}$ or simply $b$ for this
pre-action belief unless the distinction between pre-action and post-observation
beliefs is material.

For a first implementation, the operational belief can be a finite-dimensional
summary,
\begin{equation}
  \widehat b_{it}
  =
  (\widehat p^{primary}_{it},\widehat W^{out}_{it},
   \widehat p^{promo}_{it},\widehat p^{revolver}_{it},
   \widehat p^{stress}_{it},\widehat p^{attr}_{it},
   \widehat M^{lat}_{it},\widehat Rel^{lat}_{it}),
  \label{eq:operational-belief}
\end{equation}
with uncertainty and evidence grade attached to entries that materially affect
NPV. This is an implementable approximation to the belief-state object, not a
claim that the coordinates are uniquely correct.

\section{Controlled stochastic process}

\subsection{Action space by regime}

The feasible action set must depend on lifecycle regime and governance state:
\begin{equation}
  a_{it}\in\mathcal{A}_t(O_{it},b_{it};d,s,\pi^{gov}).
  \label{eq:feasible-actions}
\end{equation}
Representative stage-specific actions are:
\begin{itemize}[leftmargin=*]
  \item acquisition: suppress/contact, channel, offer family,
  \item underwriting or line assignment: approve/decline, line bucket, APR and
        fee terms, promo cell,
  \item existing-account treatment: line change, repricing, promo, retention,
        conversion,
  \item collections: no treatment, reminder, hardship, plan, escalation.
\end{itemize}

Product terms should be treated as a governed action vector,
\begin{equation}
  a_{it}
  =
  (\ell_{it},r^{pur}_{it},r^{cash}_{it},r^{BT}_{it},r^{promo}_{it},
   T^{promo}_{it},f_{it},q_{it},b_{it}^{signup},d_{it}^{disc}),
  \label{eq:action-vector}
\end{equation}
consistent with the main proposal's offer-vector representation.

\subsection{Macro block}

A full stochastic note should include an explicit macro law of motion. One
flexible specification is a linear-Gaussian latent macro regime with observed
releases:
\begin{align}
  M^{lat}_{i,t+1} &= A_M M^{lat}_{it} + \eta^M_{t+1},
    &\eta^M_{t+1}\sim\mathcal{N}(0,Q_M),
    \label{eq:macro-transition}\\
  M^{obs}_{i,t+1} &= H_M M^{lat}_{i,t+1} + \nu^M_{t+1},
    &\nu^M_{t+1}\sim\mathcal{N}(0,R_M),
    \label{eq:macro-observation}
\end{align}
with release lags handled through missing or stale components of
$M^{obs}_{i,t+1}$.

If the bank prefers deterministic stress scenarios, then
\eqref{eq:macro-transition} can be replaced by an exogenous scenario path. But
for a committee-level stochastic-process note, an explicit macro block is the
cleanest baseline because the main proposal repeatedly makes macro state a
structural driver of spend, payment, loss, and attrition
\citep{ganong2019unemployment,johnson2004rebates,gross2001liquidity,baker2020epidemic}.

\subsection{Usage, wallet, and routing block}

Let total cardable category-$j$ spending opportunity satisfy
\begin{equation}
  \log(1+C_{ij,t+1})
  =
  \mu_{ij}^{C} + \beta_j^{M\top} M^{lat}_{i,t+1}
  + \beta_j^{X\top} X_{it}^{obs}
  + \beta_j^{U\top} \Theta^{repay}_{it}
  + \varepsilon^C_{ij,t+1},
  \label{eq:cardable-spend-process}
\end{equation}
where $C_{ij,t+1}$ is latent total cardable spend opportunity in category $j$.
This lifts the proposal's category-spend specification into the latent dynamic
system.

Let routing shares across payment rails be
\begin{equation}
  q_{ij,t+1}(a)
  =
  \softmax\Bigl(
    \alpha_j + \Gamma_j O_{it} + \Lambda_j U_{it} + \Psi_j a_{it}
  \Bigr),
  \label{eq:routing-softmax}
\end{equation}
with components corresponding to new-bank card, old-bank card, outside-card,
debit/deposit, and cash/other rails. Then observed or latent rail-specific spend
is
\begin{equation}
  S^{rail}_{ij,t+1}(a) = q^{rail}_{ij,t+1}(a)\,C_{ij,t+1}.
  \label{eq:rail-specific-spend}
\end{equation}
This gives a true stochastic-process interpretation to the main proposal's
wallet accounting and spend-source decomposition.

The outside-wallet state evolves as
\begin{equation}
  W^{out}_{i,t+1}
  =
  h_W(W^{out}_{it},M^{lat}_{i,t+1},G_{it},\Theta^{route}_{it},a_{it})
  + \varepsilon^W_{i,t+1},
  \label{eq:outside-wallet-transition}
\end{equation}
and the offer-specific movable-wallet state is a derived functional
\begin{equation}
  W^{mov}_{it}(a)
  =
  m(U_{it},O_{it},a,d,s,\pi).
  \label{eq:movable-wallet-functional}
\end{equation}

\subsection{Lifecycle and activation block}

The lifecycle state transition is controlled and stage-dependent:
\begin{equation}
  \Pr(G_{i,t+1}=g'\mid G_{it}=g,O_{it},U_{it},a_{it},s)
  = p_G(g'\mid g,O_{it},U_{it},a_{it},s).
  \label{eq:lifecycle-transition}
\end{equation}
For adoption and first use, a useful two-step structure is
\begin{align}
  A_{im,t+1}^{*}
    &= \alpha_m^{\top} O_{it} + \gamma_m^{\top} U_{it} + \xi_{im,t+1},
    \qquad
    A_{im,t+1}=\1\{A^*_{im,t+1}>0\},
    \label{eq:adoption-transition}\\
  U^{use}_{ijm,t+1}
    &= \beta_m^{\top} O_{it} + \phi_m A_{im,t+1} + \rho_m^{\top} U_{it}
       + \epsilon_{ijm,t+1},
    \label{eq:conditional-use-utility}
\end{align}
with the new-bank-card rail choice generated through the routing system in
\eqref{eq:routing-softmax}. This retains the proposal's adoption/use logic while
placing it inside a monthly controlled process.

\subsection{Balance, payment, and line block}

The stock-flow identity must hold pathwise:
\begin{align}
  B_{i,t+1}^{pur}
    &= B_{it}^{pur} + S^{newbank,pur}_{i,t+1} - Pay^{pur}_{i,t+1}
       - Chg^{pur}_{i,t+1},
    \label{eq:balance-pur}\\
  B_{i,t+1}^{cash}
    &= B_{it}^{cash} + S^{cashadv}_{i,t+1} - Pay^{cash}_{i,t+1}
       - Chg^{cash}_{i,t+1},
    \label{eq:balance-cash}\\
  B_{i,t+1}^{BT}
    &= B_{it}^{BT} + BT_{i,t+1} - Pay^{BT}_{i,t+1}
       - Chg^{BT}_{i,t+1},
    \label{eq:balance-bt}\\
  B_{i,t+1}^{fee}
    &= B_{it}^{fee} + Fee_{i,t+1} - Pay^{fee}_{i,t+1},
    \label{eq:balance-fee}\\
  B_{i,t+1}^{int}
    &= B_{it}^{int} + Int_{i,t+1} - Pay^{int}_{i,t+1},
    \label{eq:balance-int}\\
  B_{i,t+1}^{tot}
    &= B_{i,t+1}^{pur}+B_{i,t+1}^{cash}+B_{i,t+1}^{BT}
       +B_{i,t+1}^{fee}+B_{i,t+1}^{int},
    \label{eq:balance-total}
\end{align}
where the flows indexed by $t+1$ are realized during $(t,t+1]$ after action
$a_{it}$ is chosen. The interest, fee, and charge-off components are generated
by the action, subledger, and delinquency states.

Payment behavior should be controlled by latent repayment state and observed
balance burden:
\begin{equation}
  Pay_{i,t+1}
  =
  g_P\bigl(B_{it}^{tot},\mathrm{Util}_{it},G_{it},D_{it},O_{it},\Theta^{repay}_{it},a_{it},M^{lat}_{i,t+1}\bigr)
  + \varepsilon^P_{i,t+1}.
  \label{eq:payment-transition}
\end{equation}

Line and term dynamics are also controlled:
\begin{equation}
  L_{i,t+1}
  =
  g_L(L_{it},B_{it}^{tot},D_{it},O_{it},a_{it},\pi^{gov}) + \varepsilon^L_{i,t+1}.
  \label{eq:line-transition}
\end{equation}

\subsection{Risk, delinquency, and recovery block}

The delinquency/default process is a controlled finite-state Markov chain or
hazard model. One flexible representation is
\begin{equation}
  \Pr(D_{i,t+1}=d'\mid D_{it}=d,G_{it},B_{it}^{tot},Pay_{i,t+1},L_{it},O_{it},U_{it},a_{it},M^{lat}_{i,t+1})
  = p_D(d'\mid \cdot).
  \label{eq:delinquency-transition}
\end{equation}

Expected exposure at default is tied to the subledgers and undrawn behavior:
\begin{equation}
  EAD_{i,t+1}
  =
  B_{i,t+1}^{pur}+B_{i,t+1}^{cash}+B_{i,t+1}^{BT}+B_{i,t+1}^{fee}+B_{i,t+1}^{int}
  + UndrawnDraw_{i,t+1}.
  \label{eq:ead-identity}
\end{equation}

LGD and recovery evolve conditionally on default state and collections action:
\begin{align}
  Rec^{work}_{i,t+1:t+W}
  &= g_R(D_{i,t+1},O_{it},U_{it},a_{it},M^{lat}_{i,t+1})
     + \varepsilon^R_{i,t+1},
    \label{eq:recovery-process}\\
  LGD_{i,t+1}
  &=
  1-\frac{\delta^{work}Rec^{work}_{i,t+1:t+W}}
          {\max(EAD_{i,t+1},\varepsilon_{EAD})}.
  \label{eq:lgd-identity}
\end{align}
Here $Rec^{work}_{i,t+1:t+W}$ is discounted recovery over the workout window,
conditional on default or charge-off. In implementation, LGD is defined only on
the default/exposure population and must reconcile to approved recovery and
charge-off conventions.

These equations make the main proposal's PD/LGD/EAD structure part of the full
state process rather than a detached risk appendix.

\subsection{Rewards, promotions, attrition, conversion, and relationship block}

Promotion state should evolve explicitly:
\begin{equation}
  Rwd_{i,t+1}
  =
  g_{Rwd}(Rwd_{it},S^{newbank}_{i,t+1},a_{it},\Theta^{promo}_{it})
  + \varepsilon^{Rwd}_{i,t+1},
  \label{eq:promo-transition}
\end{equation}
with latent promo response state
\begin{equation}
  \Theta^{promo}_{i,t+1}
  =
  A_{\Theta,p}\Theta^{promo}_{it} + \Gamma_p a_{it} + \eta^p_{i,t+1}.
  \label{eq:promo-latent-transition}
\end{equation}

Attrition and dormancy should also be explicit:
\begin{equation}
  \Pr(G_{i,t+1}\in\{\mathrm{dormant},\mathrm{closed}\}
    \mid G_{it},O_{it},U_{it},S_{it},B_{it},a_{it})
  = p_A(\cdot\mid\cdot),
  \label{eq:attrition-transition}
\end{equation}
and product conversion should satisfy
\begin{equation}
  \Pr(Product_{i,t+1}=p'\mid Product_{it}=p,G_{it},O_{it},U_{it},a_{it})
  = p_{conv}(p'\mid p,\cdot).
  \label{eq:conversion-transition}
\end{equation}

Finally, relationship spillover should have a dynamic block rather than appear
only as a terminal scalar:
\begin{equation}
  Rel^{lat}_{i,t+1}
  =
  g_{Rel}(Rel^{lat}_{it},Rel^{obs}_{it},a_{it},G_{it},M^{lat}_{i,t+1})
  + \varepsilon^{Rel}_{i,t+1}.
  \label{eq:relationship-transition}
\end{equation}

\section{Observation system}

A full state-space note must say what the bank observes and how. The observation
system is module-specific.

\subsection{Transactional and statement observations}

Observed new-bank-card and old-bank-card spend are exact or nearly exact ledger
aggregates conditional on source-system quality:
\begin{align}
  Y_{ij,t+1}^{new} &= S^{newbank}_{ij,t+1} + \nu^{new}_{ij,t+1},
    \label{eq:obs-new-spend}\\
  Y_{ij,t+1}^{old} &= S^{oldbank}_{ij,t+1} + \nu^{old}_{ij,t+1},
    \label{eq:obs-old-spend}
\end{align}
with $\nu$ typically zero or small reconciliation noise. Statement balances,
payments, fees, and interest are observed through the ledger:
\begin{align}
  Y_{i,t+1}^{B} &= B_{i,t+1}^{tot} + \nu^B_{i,t+1},
    \label{eq:obs-balance}\\
  Y_{i,t+1}^{Pay} &= Pay_{i,t+1} + \nu^{Pay}_{i,t+1}.
    \label{eq:obs-payment}
\end{align}

\subsection{Risk and bureau observations}

Delinquency and charge-off are observed discrete states,
\begin{equation}
  Y_{i,t+1}^{D} = h_D(D_{i,t+1}) + \nu^D_{i,t+1},
  \label{eq:obs-delinquency}
\end{equation}
while bureau refreshes are periodic and noisy observations of broader external
credit condition,
\begin{equation}
  Y_{i,t+1}^{bureau}
  =
  h_{bureau}(U_{i,t+1},D_{i,t+1},B_{i,t+1}^{tot}) + \nu^{bureau}_{i,t+1}.
  \label{eq:obs-bureau}
\end{equation}

\subsection{Wallet and routing anchors}

The bank does not directly observe outside wallet every month. It may observe
partial anchors such as balance transfers into the card, payoff events,
account-aggregation samples, network aggregates, or deposit-rail substitution
signals. Represent these as sparse or partial measurements:
\begin{equation}
  Y_{i,t+1}^{wallet}
  =
  h_{wallet}(W^{out}_{i,t+1},\Theta^{route}_{i,t+1},O_{i,t+1}) + \nu^{wallet}_{i,t+1},
  \label{eq:obs-wallet}
\end{equation}
where missingness itself is informative and must be retained in the filter.

\subsection{Macro observations}

Observed macro data are release-lagged measurements of the latent macro regime,
already written in \eqref{eq:macro-observation}. The note should treat revised
macro data as different objects from the as-of-decision releases used in the
filter.

\section{Bayesian filtering and belief-state recursion}

The full posterior satisfies
\begin{equation}
  b_{i,t+1}(u')
  \propto
  p(Y_{i,t+1},O_{i,t+1}\mid u',O_{it},a_{it},s,\pi^{down})
  \int p(u'\mid u,O_{it},a_{it},s,\pi^{down})b_{it}(u)\,du.
  \label{eq:generic-filter}
\end{equation}
where $O_{i,t+1}$ is the approved next decision state constructed from raw
signals. If that construction is deterministic, the likelihood includes the
state-construction rule rather than an additional random observation.
For the full enterprise model, exact filtering is generally infeasible. The
proposal should therefore define a blockwise filtering architecture.

\paragraph{Macro block.}
If \eqref{eq:macro-transition}--\eqref{eq:macro-observation} is used, the macro
posterior can be updated with Kalman-style state-space recursions
\citep{durbinkoopman2012ssm}.

\paragraph{Lifecycle and delinquency blocks.}
Finite-state transitions for $G_t$ and $D_t$ can use exact hidden Markov model
recursions conditional on observed transactional and delinquency signals.

\paragraph{Latent repayment and promo response blocks.}
If these are modeled as approximately Gaussian dynamic linear states, posterior
updates can use dynamic-linear-model recursions \citep{westharrison1997bfdm}.

\paragraph{Outside-wallet and routing block.}
Because the observation system is sparse and nonlinear, this block should be
filtered with particle or sequential Monte Carlo approximations when the bank
requires full posterior uncertainty \citep{doucet2001smc}.

The implementation may retain a structured approximate posterior,
\begin{equation}
  b_{it}(u)
  \approx
  b^{M}_{it}(M^{lat})
  b^{wallet}_{it}(W^{out},\Theta^{route})
  b^{repay}_{it}(\Theta^{repay})
  b^{promo}_{it}(\Theta^{promo})
  b^{rel}_{it}(Rel^{lat}),
  \label{eq:factorized-belief}
\end{equation}
provided the approximation error is treated as part of model risk rather than
ignored.

\begin{proposition}[Belief-state sufficiency]
Under the controlled Markov transition and observation equations above, any
history-dependent policy has an equivalent policy measurable with respect to the
current observed state $O_{it}$ and belief state $b_{it}$, with the same
finite-horizon value. The control problem is therefore an MDP on belief space
\citep{smallwood1973pomdp,kaelbling1998pomdp}.
\end{proposition}

\begin{proof}
By the tower property, any continuation value depends on past history only
through the conditional distribution of the latent state and the observed state
that indexes the transition, reward, and feasible-action correspondence. The
filter in \eqref{eq:generic-filter} constructs exactly that conditional
distribution. Replacing a history-dependent policy by one measurable with
respect to $(b_{it},O_{it})$ therefore preserves the distribution of future
rewards and the resulting value.
\end{proof}

\section{Reward and valuation mapping}

The reward must be compiled from the same economic decomposition already
proposed, not replaced by an abstract utility index. Conditional on the full
state, action, next state, and realized monthly signals,
\begin{align}
  r_{i,t+1}(X_{it},a_{it},X_{i,t+1},Y_{i,t+1};s,\pi^{down})
  &=
  \Delta \PPNR_{i,t+1}(X_{it},a_{it},X_{i,t+1};s,\pi^{down})
  \nonumber\\
  &\quad
  - \Delta \EL_{i,t+1}(X_{it},a_{it},X_{i,t+1};s,\pi^{down})
  \nonumber\\
  &\quad
  - \Delta Kchg_{i,t+1}(X_{it},a_{it},X_{i,t+1};s,\pi^{down})
  \nonumber\\
  &\quad
  - \Delta Tax_{i,t+1}(X_{it},a_{it},X_{i,t+1};s,\pi^{down})
  \nonumber\\
  &\quad
  + \Delta RelValue_{i,t+1}(X_{it},a_{it},X_{i,t+1};s,\pi^{down}).
  \label{eq:full-state-reward}
\end{align}
The belief-state reward is then
\begin{equation}
  \bar r_{it}(b_{it},O_{it},a_{it};s,\pi^{down})
  =
  \E\!\left[
    r_{i,t+1}(X_{it},a_{it},X_{i,t+1},Y_{i,t+1};s,\pi^{down})
    \mid b_{it},O_{it},a_{it},s,\pi^{down}
  \right].
  \label{eq:belief-reward}
\end{equation}

The path engine in the main proposal is then reinterpreted as simulation of the
state-space system under candidate action and baseline action, with the
subledger identities enforced exactly and the reward compiled from the simulated
paths.

\section{Constrained policy value and Bellman recursion}

Let $H$ denote the terminal valuation month for this note. The fixed-policy
value under a candidate rule $\varphi$ and downstream policy bundle
$\pi^{down}$ is
\begin{equation}
  \begin{split}
  V_t^{\varphi;\pi^{down}}(b_{it},O_{it};s)
  =
  \E^{\varphi,\pi^{down}}\!\Biggl[
    \sum_{h=t}^{H-1}\delta^{h-t}
      \bar r_h(b_{ih},O_{ih},a_{ih};s,\pi^{down})
    \\
    + \delta^{H-t}\Delta TV_{iH}(b_{iH},O_{iH};s,\pi^{down})
    \mid b_{it},O_{it}
  \Biggr].
  \end{split}
  \label{eq:policy-value}
\end{equation}

The constrained Bellman recursion is
\begin{align}
  V_t^{\star}(b,O;s)
  &=
  \max_{a\in\mathcal{A}_t(b,O;d,s,\pi^{gov})}
  \Bigl\{
    \bar r_t(b,O,a;s,\pi^{down})
    \nonumber\\
    &\qquad\qquad
    + \delta\,\E\left[V_{t+1}^{\star}(b',O';s)\mid b,O,a,s,\pi^{down}\right]
  \Bigr\}
  \label{eq:bellman-main}
\end{align}
with pointwise hard admissibility
\begin{equation}
  a\in \mathcal{A}^{hard}_{t}(O,b,d,s,\pi^{gov}),
  \label{eq:hard-feasible}
\end{equation}
and cumulative constraint costs
\begin{equation}
  J^{\varphi}_{c,k}(b,O;s)
  =
  \E^{\varphi,\pi^{down}}\!\left[
    \sum_{h=t}^{H-1}\delta^{h-t}
    c_{k,h}(b_{ih},O_{ih},a_{ih};s)
    \mid b_{it}=b,O_{it}=O
  \right]
  \le \kappa_k.
  \label{eq:cumulative-constraints}
\end{equation}
Loss, capital, liquidity, operational, complaint, and fairness constraints are
usually cumulative or portfolio-level objects, not merely one-period
inequalities. They can be handled through remaining-budget state augmentation or
Lagrangian relaxations when the policy class and evidence support optimization.

A fully observed MDP is the special case in which $U_{it}$ is either empty or
exactly observed and $b_{it}$ degenerates to a point mass.

\section{Module-to-process mapping}

\begin{table}[t]
\centering
\caption{How the existing proposal modules enter the full stochastic process.}
\label{tab:module-map}
\begin{tabular}{L{0.18\linewidth}L{0.21\linewidth}L{0.20\linewidth}L{0.28\linewidth}}
\toprule
Module & Transition role & Observation role & Reward / constraint role \\
\midrule
Macro and labor state
  & latent and observed macro regime transitions
  & release-lagged macro signals
  & scenario conditioning; stress and drift constraints \\
Adoption / activation / use
  & lifecycle transitions and route choice
  & application, booking, activation, first use, repeat use
  & gates future value and contact-policy feasibility \\
Wallet and cannibalization
  & outside wallet and routing dynamics
  & partial anchors from bank book, bureau, network, or experiments
  & determines all-bank incremental spend versus raw new-card spend \\
Balances / payments / lines
  & subledger and payment transitions
  & statement balances, payments, interest, fees
  & finance charges, funding, EAD, line constraints \\
Delinquency / losses / recoveries
  & delinquency, default, cure, and recovery transitions
  & delinquency buckets, charge-off, recoveries, bureau refreshes
  & expected loss, capital, downside-risk constraints \\
Rewards / promos
  & promo-state and latent response evolution
  & qualification, redemption, reward liability, post-promo decay
  & reward cost and temporary-versus-durable lift \\
Attrition / conversion / relationship
  & dormancy, closure, conversion, retention, relationship state
  & activity, closure, other-product observations
  & horizon length, relationship value, retention constraints \\
\bottomrule
\end{tabular}
\end{table}

\section{Identification, estimation, and evidence discipline}

The stochastic-process note should classify each equation or block as one of the
following:
\begin{enumerate}[label=\arabic*.]
  \item accounting identity,
  \item probabilistic factorization,
  \item literature-supported control or filtering principle,
  \item bank-specific empirical object,
  \item governance or admissibility convention.
\end{enumerate}

The literature directly supports:
\begin{itemize}[leftmargin=*]
  \item the MDP / dynamic treatment regime framing
        \citep{puterman1994mdp,murphy2003dynamic};
  \item a credit-card precedent for sequential NPV optimization
        \citep{trench2003bankone};
  \item dynamic stochastic repayment and collections processes
        \citep{chehrazi2015dynamicvaluation,chehrazi2019collections};
  \item belief-state or POMDP control language
        \citep{smallwood1973pomdp,kaelbling1998pomdp};
  \item Bayesian filtering language
        \citep{durbinkoopman2012ssm,westharrison1997bfdm,doucet2001smc};
  \item offline evaluation and safe-improvement tools
        \citep{dudik2011dr,jiangli2016dr,thomasbrunskill2016ope,
        huwager2021pomdpope,laroche2019spibb,levine2020offline,
        fujimoto2019bcq,kumar2020cql,khraishiokhrati2022offlinecredit}.
\end{itemize}

The literature does not identify this bank's:
\begin{itemize}[leftmargin=*]
  \item exact latent state vector,
  \item outside-wallet and movable-wallet laws of motion,
  \item reward ledger detail,
  \item feasible action correspondence,
  \item behavior-policy logging quality,
  \item or policy-promotion threshold.
\end{itemize}

Those objects must be justified with internal experiments, quasi-experiments,
logged-action support, finance/risk reconciliation, and governance approval.

\section{Implementation and governance ladder}

The stochastic process implies four component interfaces: \code{simulate_policy}
for fixed-policy forward simulation, \code{score_action_grid} for finite-menu
line/offer/retention decisions, \code{evaluate_policy_offline} for comparing a
candidate rule $\varphi$ against the logged policy $\mu$, and
\code{register_policy} for versioning policy rules and their allowed-use scope.
These are NPV-component interfaces, not ownership of the bank's campaign,
underwriting, line-management, or servicing platforms.

The correct rollout sequence is:
\begin{enumerate}[label=\arabic*.]
  \item build the fixed-policy forward simulator under the incumbent or declared
        policy bundle;
  \item calibrate and reconcile transitions, cash-flow compiler, and risk
        outputs;
  \item evaluate challenger policies offline with support diagnostics, behavior-
        policy probabilities $\mu(a_t\mid H_t)$ where available, and policy-value
        estimators \citep{dudik2011dr,jiangli2016dr,thomasbrunskill2016ope};
  \item allow constrained challenger optimization only for narrow approved
        subproblems, preferably with discrete action grids and conservative
        support-aware methods \citep{ernst2005fqi,laroche2019spibb,
        fujimoto2019bcq,kumar2020cql};
  \item allow live influence only through governed experimentation and rollback
        controls.
\end{enumerate}

This ordering is not optional. Without it, the belief-state control language can
launder weakly identified modules into apparently precise policy recommendations.

\section{Source-support audit}
\label{sec:audit}

This note is a mathematical assembly of the current NPV proposal into a full
state-space control specification. It is not a completed empirical validation
package. Table~\ref{tab:audit-ledger} records the support class for the main
claims.

\begin{longtable}{L{0.25\linewidth}L{0.23\linewidth}L{0.42\linewidth}}
\caption{Claim-support ledger for the full stochastic-process note.}
\label{tab:audit-ledger}\\
\toprule
Design claim & Support class & What remains bank-specific \\
\midrule
\endfirsthead
\toprule
Design claim & Support class & What remains bank-specific \\
\midrule
\endhead
Credit-card customer value should be modeled as a sequential control problem.
  & Direct literature support plus card-management precedent
  & Exact state/action granularity and policy scope \\
Partial observability requires a belief-state or POMDP formulation.
  & Direct control/filtering support; proposal derivation for wallet and latent
    response states
  & Exact posterior approximation and state coordinates \\
Macro, spend, payment, loss, reward, and attrition blocks must be assembled in
one stochastic system.
  & Proposal derivation from reused equations and literature-supported
    mechanisms
  & Exact parametrization and empirical calibration \\
Policy should condition on the posterior distribution, not only a point estimate.
  & Direct filtering and POMDP support
  & Chosen finite-dimensional approximation and evidence threshold \\
Offline policy evaluation must precede challenger optimization.
  & Direct OPE / safe-improvement literature support
  & Logged behavior-policy quality, overlap, support, and promotion rule \\
Safe live rollout is not implied by the literature alone.
  & Governance conclusion from the limits of the cited methods
  & All live thresholds, rollback rules, and compliance approval \\
\bottomrule
\end{longtable}

\section{Conclusion}

A committee-ready dynamic note for the credit-card NPV component cannot stop at
a high-level POMDP backbone. It must specify the whole stochastic process:
macro block, usage and routing, wallet and cannibalization, balances, payments,
delinquency and recovery, rewards and promotions, attrition and conversion,
relationship state, observation equations, posterior update, reward mapping, and
constrained policy recursion. This note does that by reassembling the current
proposal's modules into one belief-state control model.

The resulting object is mathematically strong enough to define forward
simulation under a fixed policy and conceptually strong enough to support later
offline policy evaluation and constrained challenger optimization. It remains
properly conservative: the literature supports the control and filtering
language, but the empirical state content, parameterization, and rollout
thresholds remain issuer-specific obligations.

\appendix

\section{Notation summary}

\begin{longtable}{L{0.18\linewidth}L{0.72\linewidth}}
\toprule
Symbol & Meaning \\
\midrule
\endfirsthead
\toprule
Symbol & Meaning \\
\midrule
\endhead
$i$ & customer / account / prospect / application unit \\
$t$ & monthly time index \\
$d$ & decision context \\
$a, a_0(d)$ & action and baseline action \\
$s$ & scenario and valuation-semantics bundle \\
$\pi^{down}$ & downstream operating policy bundle assumed in valuation \\
$\mu$ & historical logged behavior policy that generated observed actions \\
$\varphi$ & candidate or challenger decision rule \\
$\pi^{gov}$ & governance and hard-admissibility bundle \\
$O_t$ & observed decision state \\
$U_t$ & latent state \\
$X_t$ & full state $(O_t,U_t)$ \\
$b_t$ & posterior belief over the latent state \\
$G_t$ & lifecycle regime \\
$M_t^{lat},M_t^{obs}$ & latent and observed macro state \\
$W_t^{out},W_t^{mov}$ & outside wallet and movable wallet \\
$B_t$ & balance-subledger vector \\
$D_t$ & delinquency/default/cure state \\
$Rwd_t$ & reward / promotion state \\
$Rel_t$ & relationship state \\
$\Delta CF_t$ & one-period incremental cash-flow primitive \\
$\Delta \NPV$ & discounted incremental value \\
\bottomrule
\end{longtable}

\section{Assumptions ledger}

\begin{longtable}{L{0.22\linewidth}L{0.68\linewidth}}
\toprule
Assumption & Role in the model \\
\midrule
\endfirsthead
\toprule
Assumption & Role in the model \\
\midrule
\endhead
Controlled Markov structure & Defines the latent-state law of motion under action \\
Observation model & Links latent/full state to observed bank signals \\
Belief-state sufficiency & Justifies control on posterior state rather than full history \\
Subledger accounting & Prevents impossible balance / revenue / loss paths \\
Constraint admissibility & Restricts optimization to compliant and approved actions \\
Offline evaluation assumptions & Required for policy-value estimation from logs \\
Evidence-grade separation & Prevents weak latent modules from being silently promoted to causal policy value \\
\bottomrule
\end{longtable}

\bibliographystyle{plainnat}
\bibliography{credit_card_npv_component_refs}

\end{document}
